\chapter{绪论}
\label{cha:introduction}

\section{研究背景与意义}
\label{sec:introduction-background}

随着智能网联汽车、路侧单元与云平台之间的协同程度不断提高，车联网已经从单纯的车载通信场景扩展为融合感知、决策、控制与服务分发的复杂网络系统。车辆与车辆、车辆与路侧基础设施之间持续交换位置、速度、诊断状态与控制指令等数据，这一方面提高了交通效率与行车安全，另一方面也显著扩大了攻击面。一旦通信链路、车载终端或边缘节点遭受伪造报文、重放、拒绝服务或异常控制指令注入等攻击，轻则导致业务异常，重则可能引发交通安全风险。因此，在车联网环境中构建及时、可靠且具备工程可落地性的入侵检测机制，已经成为智能交通安全研究中的重要问题之一\cite{placeholder_iov_security_survey,placeholder_vanet_ids_survey}。

传统入侵检测系统通常依赖集中式日志汇聚与统一训练流程，即将多源数据收集到中心节点后进行特征工程、模型训练与全局判决。这类方法在静态或可控网络环境中较易实施，但在车联网场景下往往面临三方面限制。首先，车辆节点的算力、能耗与链路状态存在显著差异，持续上传原始数据会带来较高的带宽成本与调度压力。其次，车载数据天然包含位置轨迹、行为模式等敏感信息，直接集中存储存在隐私泄露风险。最后，不同车辆或不同道路环境下的数据分布并不一致，简单套用中心化训练策略往往难以兼顾泛化能力、通信代价与部署效率\cite{placeholder_distributed_ids,placeholder_privacy_preserving_ids}。

联邦学习为上述问题提供了新的解决思路。该方法强调“数据不出本地”，各参与节点仅上传模型更新或统计信息，由中心协调端完成聚合，从而在一定程度上兼顾数据可用性与隐私保护需求\cite{placeholder_federated_learning_survey,placeholder_fedavg}。然而，联邦学习并不天然适用于车联网入侵检测。一方面，车辆节点参与训练具有动态性与异质性，若所有节点在每轮中无差别参与，将导致训练时延和通信量迅速升高；另一方面，模型更新本身也会带来不小的传输负担，若缺少压缩和筛选机制，则系统整体效率仍然受限。基于此，本文围绕联邦学习与车联网入侵检测的结合问题，研究一个面向工程实现的原型系统，并在真实代码基础上讨论其设计、实现与实验验证过程。

从工程应用角度看，本研究具有以下意义：其一，为车联网入侵检测任务提供了一套从数据预处理、联邦训练到本地评估的完整链路，有助于将算法验证与系统实现结合起来；其二，通过引入节点选择、Top-K 稀疏化与定点量化等策略，探索在资源受限网络中降低通信开销的可行路径；其三，将配置快照、日志、运行元数据与实验报告统一沉淀到输出目录中，为后续复现实验、撰写论文和补充正式数据集验证提供了证据基础。由于当前开题报告与任务书正文尚未完成结构化提取，相关原始措辞与文献脉络在本文中暂以受控占位形式保留，后续可结合正式材料进一步细化\cite{placeholder_proposal_taskbook}。

\section{国内外研究现状}
\label{sec:introduction-related-work}

围绕车联网安全问题，国内外研究已从早期基于规则与签名的检测方法，逐步发展到结合统计学习和深度学习的智能化检测方法。规则驱动方法实现简单、解释性较强，适用于已知攻击模式的快速识别，但对未知攻击与复杂异常行为的适应能力有限。随着车载总线日志、网络流量与多源传感数据逐渐丰富，研究者开始采用支持向量机、随机森林、卷积神经网络、循环神经网络等模型对异常流量与攻击事件进行识别，以期提升复杂场景下的检测精度与泛化能力\cite{placeholder_ml_ids_survey,placeholder_deep_ids_survey}。不过，这类方法多数默认数据可以集中收集，较少充分考虑车联网环境中持续在线训练、链路受限与隐私保护并存的现实约束。

在联邦学习方向，已有研究证明该范式能够在医疗、金融、移动终端与边缘计算等场景中支持跨节点协同建模，其中 FedAvg 是最具代表性的基础聚合方法之一，其思想是各客户端完成若干轮本地更新后，由服务器对模型参数进行加权平均\cite{placeholder_fedavg,placeholder_federated_learning_survey}。随后，FedProx 等方法通过在本地目标函数中加入近端正则项，缓解非独立同分布数据和系统异质性带来的训练不稳定问题\cite{placeholder_fedprox}。这些方法为车联网入侵检测提供了理论基础，但直接迁移到车辆场景时仍面临节点数量波动大、训练参与机会不均衡、更新上传成本高等挑战。

近年来，也有研究尝试将联邦学习用于工业互联网、物联网和车联网的异常检测任务，通过分布式训练方式提升数据利用率并降低原始数据外泄风险\cite{placeholder_iov_fl_ids,placeholder_edge_security_fl}。然而，已有方案通常聚焦于联邦学习框架本身或单一模型效果，对工程侧的完整链路描述相对不足，例如缺少标准化的数据预处理流程、配置管理机制、训练证据留存方式以及通信统计分析。同时，一些工作虽然关注压缩通信或客户端选择，但往往分别讨论节点筛选与模型压缩，缺少面向统一系统的综合实现与对照验证。基于这一现状，本文将研究重点放在“可运行、可复现、可分析”的车联网联邦入侵检测原型构建上，并在实验部分讨论改进策略相对标准 FedAvg 基线的收益与局限。

\section{研究内容与目标}
\label{sec:introduction-objective}

本文以当前仓库中的 Python 项目和华南理工大学论文模板为基础，面向“基于联邦学习的车联网入侵检测系统设计与实现”这一题目组织论文内容。结合现有代码与实验材料，本文的研究内容主要包括以下几个方面。

第一，构建完整的系统运行链路。项目在 `main.py` 中提供统一命令行入口，支持 `preprocess`、`train` 与 `test` 三种模式；其中，预处理阶段负责读取本地表格数据并生成标准化训练集、验证集和标签映射，训练阶段负责联邦轮次调度、客户端本地更新与全局聚合，测试阶段则加载全局模型并对验证集进行评估。围绕这一主流程，系统还引入了配置加载、随机种子设置、日志记录与运行元数据保存等支撑模块，使得实验过程具备较好的复现性与可追踪性。

第二，设计适用于车联网场景的联邦训练改进策略。现有实现以 FedAvg/FedProx 框架为基础，在客户端选择阶段综合考虑节点算力、电量与信道质量，并通过轻量公平性惩罚减弱同一节点被频繁选中的问题；在模型更新上传阶段，对参数增量执行 Top-K 稀疏化与定点量化，以降低通信量并形成可度量的压缩收益。上述策略并非脱离代码的理论设想，而是已经在 \texttt{federated\_learning.py} 中落成具体实现，并可通过训练历史与报告文件进行核验。

第三，构建轻量化的入侵检测模型与评价体系。项目使用 `models.py` 中定义的 `LightweightCNNLSTM` 作为检测器，通过一维卷积提取局部特征，再由 LSTM 建模序列关系，最终输出攻击类别预测。为全面分析系统效果，本文不仅关注准确率、召回率、精确率和 F1 值等检测指标，还结合通信量、训练耗时、隐私代理分数等工程指标，对改进方案与基线方案进行综合比较。

基于以上内容，本文的研究目标可以概括为：在不上传原始车联网数据的约束下，设计并实现一套具备数据预处理、联邦训练、本地评估与证据报告能力的入侵检测系统；在现有演示实验条件下验证动态节点选择与参数压缩机制对通信成本和训练效率的改善作用；同时保持对系统边界的审慎表述，不将当前小样本结果直接推广为大规模真实场景的最终结论。对于开题报告、任务书中尚待补充的原始表述，本文先按已验证事实给出概括性描述，后续再根据正式材料补全文字与文献来源。

\section{论文结构安排}
\label{sec:introduction-structure}

本文正文共分为五章。第一章为绪论，介绍课题背景、研究意义、国内外研究现状以及本文的研究内容与目标。第二章为相关技术与理论基础，重点阐述车联网入侵检测、联邦学习、FedAvg/FedProx、轻量化 CNN-LSTM 模型、节点选择机制、参数压缩与评价指标等基础内容。第三章为系统分析与总体设计，结合项目实际代码结构说明系统目标、使用场景、模块划分、数据流和运行流程。第四章为系统实现，围绕数据预处理、客户端切分、联邦训练、模型聚合、通信统计与结果落盘等关键实现细节展开说明。第五章为实验设计与结果分析，给出实验目标、参数设置、基线方案、结果对比和局限性讨论。正文之后为结论与展望、参考文献、致谢及附录，用于总结全文、保留佐证材料并支持后续继续完善。
